\session{10}{11/7}{1限}{3}{予測と分析の考え方}

\goals{
  \item 相関と因果の違いを、交絡という概念を使って説明できる。
  \item 予測モデル（機械学習）の基本的な考え方と、精度の評価方法を理解する。
  \item AIが出した分析結果の前提と限界を確かめる手順を持つ。
}

\guidesec{解説}

\topic{相関と因果}
2つの変数が一緒に増減する関係を\textbf{相関}と呼びます。相関があっても、一方が他方の原因であるとは限りません。
相関が見えるのに因果がない主な理由は次の3つです。
\begin{itemize}
  \item \textbf{交絡}　両方に影響する第三の要因がある。アイスの売上と水難事故はどちらも気温（夏）が原因で増える。
  \item \textbf{逆の因果}　原因と結果が逆。「売上の高い店ほど広告費が多い」は、広告が売上を生んだのではなく、売上のある店が広告に使えるのかもしれない。
  \item \textbf{偶然}　データが少ないときや、多数の変数を調べたときに、偶然の一致が見つかる。
\end{itemize}
経営データでは、「広告費と売上」「研修時間と業績」「AI導入と生産性」など、因果を主張したくなる相関が多くあります。
因果を確かめるには、比較対象（AIを導入しなかった同じような部門）を設けた実験（A/Bテスト）や、交絡を統計的に調整する分析が必要です。
AIに「因果関係はありますか」と聞いても、データからだけでは答えられません。

\topic{予測モデルの考え方}
機械学習は、データからパターンを学習し、新しいデータについて予測や判断を行う技術です。大きく3つに分かれます。
\par\noindent\begin{gtabh}{colspec={Q[l,wd=22mm] X[l] X[l]}}
  種類 & 内容 & 経営での例 \\
  教師あり学習 & 正解ラベル付きのデータから、入力と出力の関係を学習する。出力が数値なら回帰、分類なら分類。 & 需要予測（回帰）、解約しそうな顧客の判定（分類）、不良品の検知（分類） \\
  教師なし学習 & ラベルのないデータから、構造やグループを見つける。 & 顧客のセグメント分け（クラスタリング）、異常検知 \\
  強化学習 & 行動の結果として得られる報酬を最大化する方策を学習する。 & 広告配信の最適化、在庫の動的な制御 \\
\end{gtabh}
生成AI（LLM）は、教師あり学習の一種である「次のトークンの予測」を大規模に行ったものです。
経営の現場で「AIで予測する」と言うとき、多くは教師あり学習を指します。

\topic{精度の評価}
予測モデルの実力は、学習に使ったデータではなく、学習に使っていない\textbf{テストデータ}で評価します。
学習データでの精度だけを見ると、データの細部や偶然のパターンまで覚え込んで、未知のデータには当たらない\textbf{過学習}を見抜けません。
データを分割して評価を繰り返す交差検証も広く使われます。

評価の指標は、問題の性質で選びます。
\begin{itemize}
  \item \textbf{回帰}　予測値と実際の値の差（誤差）の平均。MAE（平均絶対誤差）、RMSE（二乗平均平方根誤差）など。
  \item \textbf{分類}　正解率のほか、「陽性と予測したもののうち本当に陽性の割合」（適合率）、「本当に陽性のもののうち陽性と予測できた割合」（再現率）。
        不良品の見逃し（偽陰性）と、良品を不良と判定する誤り（偽陽性）では、費用が異なるため、どちらを重視するかを業務で決める。
\end{itemize}
また、「予測モデルは単純な方法（前年同月と同じと予測する、平均を予測する）よりどれだけ良いか」という\textbf{ベースラインとの比較}が重要です。
複雑なモデルが単純な方法に勝てないことは珍しくありません。

\topic{AIの分析結果を検証する}
生成AIにデータを渡して分析させると、もっともらしい結果と解釈が返ってきます。使う前に次を確認します。
\begin{checks}
  \item 使ったデータの範囲（期間、対象、除外したもの）は何か。
  \item 置いた仮定は何か（季節性はあるか、傾向は続くか、欠損はどう扱ったか）。
  \item 精度はどう評価したか。テストデータで評価したか。ベースラインと比べたか。
  \item 結果は再現できるか。同じデータで同じ手順を繰り返して同じ結果になるか。
  \item 相関を因果として語っていないか。
  \item 予測が外れたときの影響と、外れた場合の対応は決まっているか。
\end{checks}
AIに「この分析の前提と限界を説明してください」と尋ねることは有効ですが、AIの説明もまた検証の対象です。

\topic{キーワード}
\par\noindent\begin{keywords}{}
  相関と因果 & 一緒に増減する関係と、一方が他方を引き起こす関係。相関は因果を意味しない。 \\
  交絡 & 2つの変数の両方に影響する第三の要因。見かけの相関を生む。 \\
  教師あり／教師なし／強化学習 & 機械学習の3つの種類。正解付きデータ、正解なしデータ、報酬から学習する。 \\
  過学習 & 学習データに合わせすぎて、新しいデータへの精度が下がること。 \\
  テストデータ & 精度の評価のために、学習に使わずに取っておくデータ。 \\
  適合率と再現率 & 分類の評価指標。誤りの種類（偽陽性・偽陰性）の費用に応じて重視する方を決める。 \\
\end{keywords}

\guidesec{演習課題}

\exercise{1}{AI演習：AIに簡単な予測分析を行わせ、前提と限界を説明させて妥当性を確かめる}
\instr{第9回の商品別売上の表（または自分で用意した時系列データ）を使い、次のプロンプトで2027年の予測を行わせてください。出力された予測値と説明について、解説の「AIの分析結果を検証する」のチェック項目を1つずつ当てはめ、採用できるか判断してください。}
\begin{promptbox}[予測と前提の説明]
次の表は、ある会社の商品別の年間売上（百万円）です。(1) 商品A・B・Cそれぞれの2027年の売上を予測し、(2) 使った方法と置いた仮定を説明し、(3) この予測の限界と、予測を外す可能性のある要因を挙げ、(4) 「前年と同じ」と予測する単純な方法と比べて、あなたの方法がどれだけ優れているかを、過去のデータで検証して示してください。計算はコードで行ってください。
（表）
\end{promptbox}

\exercise{2}{相関と因果を見分ける}
\instr{次の3つの主張について、「相関から因果を主張している」かどうかを判定し、交絡・逆の因果・偶然のどれが疑われるかと、因果を確かめるために必要なデータや方法を書いてください。}
\begin{itemize}
  \item (A) 「生成AIを使っている社員ほど評価が高い。だからAIを使わせれば評価が上がる」
  \item (B) 「SNS の投稿数が多い月は売上が高い。だから投稿を増やせば売上が増える」
  \item (C) 「研修を受けた店舗は受けていない店舗より売上の伸びが大きい。だから研修は効果がある」
\end{itemize}

\exercise{3}{評価指標を選ぶ}
\instr{次の2つの業務で分類モデルを使うとき、適合率と再現率のどちらを重視すべきか、偽陽性・偽陰性の費用を具体的に挙げて説明してください。}
\begin{itemize}
  \item (A) クレジットカードの不正利用の検知（不正と判定した取引は一時停止し、顧客に確認する）
  \item (B) 見込み客への営業電話の対象者の選定（「買いそう」と判定した顧客に電話する）
\end{itemize}

\guidesec{模範解答}

\begin{answer}{演習1}
AIの出力の例：「線形回帰により、2027年は商品A 513、商品B 285、商品C 407 と予測。仮定は過去の傾向が続くこと。
限界は、データが5点と少ないこと、価格や競合の情報がないこと。『前年と同じ』と比べると、過去4年の平均絶対誤差は線形回帰 12 に対し前年同値 28 で、線形回帰の方が良い。」

検証：
\begin{itemize}
  \item データの範囲：5年分の年次データのみ。月次や顧客別の情報は使っていない。
  \item 仮定：直線的な傾向が続く。商品Cの伸びが鈍化する可能性や、商品Bが下げ止まる可能性は考慮されていない。
  \item 精度の評価：「過去4年で検証」とあるが、5点のデータでの検証は信頼性が低い。AIの示した誤差の計算をコードで確かめる。
  \item 再現性：コードが示されていれば再実行できる。示されていなければ求める。
  \item 相関と因果：予測の説明に「商品Cが商品Bを奪っている」という因果の表現があれば、それはデータからは言えない。
  \item 外れた場合：2027年の商品Bが予測を上回っても下回っても、どう対応するかは別途決める必要がある。
\end{itemize}
判断：予測値は「傾向が続いた場合の目安」としては使えるが、仕入れや人員の計画の根拠にするには、月次データでの再分析と、
価格・競合の情報の追加が必要。
\end{answer}

\begin{answer}{演習2}
\begin{itemize}
  \item (A) 因果を主張している。逆の因果（評価の高い社員は新しい道具を試す意欲が高い）と交絡（能力や意欲が、AI利用と評価の両方に影響）が疑われる。
        確かめるには、同じような社員を無作為にAI利用群と非利用群に分けて評価を比べる（実験）か、過去の評価を調整した分析が必要。
  \item (B) 因果を主張している。交絡（繁忙期は投稿も売上も増える：季節性）と逆の因果（売れている月は投稿のネタが多い）が疑われる。
        確かめるには、季節を揃えたうえで投稿数を意図的に変える期間を設けて比較する。
  \item (C) 因果を主張している。交絡（意欲の高い店長の店ほど研修に参加し、売上も伸びる：自己選択）が疑われる。
        確かめるには、研修を受ける店舗を無作為に選ぶか、研修前の売上の伸びを調整して比較する。
\end{itemize}
共通する点：「AIを使う」「投稿する」「研修を受ける」のは、何らかの特徴を持つ人や店舗が自ら選んでおり、その特徴が結果にも影響する（自己選択）。
\end{answer}

\begin{answer}{演習3}
\begin{itemize}
  \item (A) 再現率を重視。偽陰性（不正を見逃す）の費用は不正利用の損失と信用の失墜で大きい。偽陽性（正常な取引を止める）の費用は顧客への確認の手間と不便で、相対的に小さい。
        ただし偽陽性が多すぎると顧客が離れるため、確認の方法（アプリでの即時確認）で費用を下げる工夫が必要。
  \item (B) 適合率を重視。偽陽性（買わない人に電話する）の費用は営業の時間と、顧客に迷惑をかける評判の低下。偽陰性（買いそうな人に電話しない）の費用は機会損失だが、
        営業の時間が限られているなら、電話した人の成約率（適合率）を上げる方が効率的。
\end{itemize}
どちらを重視するかは、モデルの性能ではなく業務の費用構造で決まる、という点が要点です。
\end{answer}

\guidesec{出席確認クイズの解説}
講義サイトの出席確認ページ（第10回）の5問です。
% 出席確認クイズ 第10回　予測と分析の考え方
%   attendance/attendance10.html から生成（解説は本ガイド用に加筆）。

\quizq{1}{「相関があっても因果があるとは限らない」ことを示す例として適切なものはどれでしょうか?}%
  {因果関係があれば相関は必ずゼロになる}%
  {\correct{アイスの売上と水難事故が共に夏に増える(気温という別の要因がある)}}%
  {相関と因果はまったく同じ意味である}%
  {相関があれば必ず因果関係がある}%
  {正解は (b)。アイスの売上と水難事故はどちらも夏に増えるため相関がありますが、一方が他方の原因ではなく、気温という共通の要因（交絡因子）があるだけです。(a)(c)(d) は相関と因果の関係を誤解しています。経営データでも「広告費と売上」のように第三の要因（季節、景気）を疑う習慣が必要です。}

\quizq{2}{「教師あり学習」の説明として適切なものはどれでしょうか?}%
  {人がルールをすべて書く}%
  {ラベルのないデータをグループ分けする}%
  {\correct{正解ラベル付きのデータから入力と出力の関係を学習する}}%
  {報酬を最大化する行動を学ぶ}%
  {正解は (c)。教師あり学習は、正解ラベル付きのデータを使って、入力から出力を予測するモデルを学習する方法です（回帰・分類）。(b) は教師なし学習（クラスタリングなど）、(d) は強化学習、(a) はルールベースの説明です。}

\quizq{3}{「過学習(オーバーフィッティング)」の説明として適切なものはどれでしょうか?}%
  {データが多すぎて計算できない}%
  {学習が足りず精度が低い}%
  {\correct{学習データに合わせすぎて、新しいデータに対する精度が下がる}}%
  {モデルが小さすぎる}%
  {正解は (c)。過学習（オーバーフィッティング）は、学習データの細部や偶然のパターンまで覚え込んでしまい、未知のデータへの予測精度が落ちる現象です。(b) は逆に学習が足りない未学習（アンダーフィッティング）の説明です。}

\quizq{4}{予測モデルの精度を適切に評価する方法はどれでしょうか?}%
  {\correct{学習に使っていないデータ(テストデータ)で確かめる}}%
  {学習データでの精度だけを見る}%
  {精度は評価しない}%
  {モデルの大きさで判断する}%
  {正解は (a)。モデルの実力は、学習に使っていないテストデータでの精度で評価します。(b) 学習データでの精度だけを見ると過学習を見抜けません。交差検証（データを分割して評価を繰り返す方法）も広く使われます。}

\quizq{5}{AIの分析結果を経営で使う前に確認すべきこととして適切なものはどれでしょうか?}%
  {\correct{前提(使ったデータ・期間・仮定)と限界}}%
  {使ったパソコンの色}%
  {出力の文字数}%
  {AIが出したので確認は不要}%
  {正解は (a)。分析結果は、使用したデータ・期間・置かれた仮定によって変わるため、前提と限界を確認してから経営判断に使います。(d) AIが出した結果も検証が必要です。(b)(c) は分析の妥当性と無関係です。}

